\documentclass[11pt,a4paper]{article}
\usepackage[margin=24mm]{geometry}
\usepackage{amsmath,amssymb,booktabs,hyperref,siunitx}
\title{SST Topology--Projection--Phase v0.2.0\\Scientific Falsifier Protocol}
\author{SST-Workbench Research Module}
\date{16 October 2026}
\newcommand{\vswirl}{\mathbf{v}_{\!\boldsymbol{\circlearrowleft}}}
\newcommand{\rhoF}{\rho_{\!f}}
\begin{document}\maketitle
\begin{abstract}
A preregistered blind geometric discrimination between 3D topological changes,
2D projection tangency, and phase-field zeros, extended by a real-source
KnotPlot/PKLSA $3_1$ trefoil Euler-pilot. Geometry qualification does not imply
particle physics, ideal-Euler reconnection or model-scale calibration.
\end{abstract}
\section{Separation of phenomena}
For optical phase defects, the topological integer is
\begin{equation}n=(2\pi)^{-1}\oint_C\nabla\phi\cdot d\boldsymbol{\ell}\in\mathbb{Z},\end{equation}
whereas a material vortex loop has circulatory transport
\begin{equation}\Gamma=\oint_C\mathbf{u}\cdot d\boldsymbol{\ell},\quad
 \boldsymbol{\omega}=\nabla\times\mathbf{u}.\end{equation}
Optical singularity motion is not automatically fluid-particle motion.
The existing encrypted DS01--DS03 synthetic test inputs are reused byte for byte:
this is regression qualification, not new independent experimental evidence.
\section{Knot and link tests}
For a polygonal closed oriented curve, three generic projections are used to
form Wirtinger--Fox integer crossing matrices $M$, with row coefficients
$2,-1,-1$. The Alexander determinant at $t=-1$ is
\begin{equation}\det(K)=|\det M'|,\end{equation}
where one row and column are deleted. The matrix determinant uses exact integer
Bareiss elimination. The control signatures are
$\det(3_1)=3$, $\det(4_1)=5$, $\det(0_1)=1$.
A determinant of unity does not prove that the curve is unknotted.
For two disjoint filament loops the Gauss integral is
\begin{equation}
\mathrm{Lk}(C_1,C_2)=\frac{1}{4\pi}\oint_{C_1}\oint_{C_2}
\frac{(d\mathbf{r}_1\times d\mathbf{r}_2)\cdot(\mathbf{r}_1-\mathbf{r}_2)}
 {|\mathbf{r}_1-\mathbf{r}_2|^3}.
\end{equation}
Tensor Gauss--Legendre quadratures of order 3 and 5 are compared. The
pre-registered acceptance bounds are $|\mathrm{Lk}_5-\operatorname{round}\mathrm{Lk}_5|<0.075$
and $|\mathrm{Lk}_5-\mathrm{Lk}_3|<0.075$ subject to centerline separation.
A valid Hopf-link control yields signed $\mathrm{Lk}=-1$ for the chosen orientation.
\section{Inviscid Euler and numerical pilot}
\begin{align}
\nabla\cdot\mathbf{u}&=0,\\
\partial_t\boldsymbol{\omega}
 +(\mathbf{u}\cdot\nabla)\boldsymbol{\omega}
 -(\boldsymbol{\omega}\cdot\nabla)\mathbf{u}&=0.
\end{align}
For smooth ideal Euler evolution, vorticity lines are advected materially and
Kelvin circulation is preserved on each material contour. True vortex
reconnection cannot be concluded from a change of 2D section counts, aggregate
circulation, or synthetic polygon surgery.
The exploratory periodic Fourier solver uses
\begin{equation}
\partial_t\hat{\mathbf u}(\mathbf{k})=
\left(I-\frac{\mathbf{k}\mathbf{k}^\mathsf{T}}{|\mathbf{k}|^2}\right)
\mathcal{F}[\mathbf{u}\times(\nabla\times\mathbf{u})],\qquad\mathbf k\neq0.
\end{equation}
It employs 2/3-dealiasing and RK4, with trilinearly interpolated material
marker velocities. The initial real-source N300 KnotPlot trefoil is normalized
into a $2\pi$-periodic dimensionless cube, and its tangent vorticity is
Gaussian-smoothed with kernel width $\sigma=0.50$ model-length units.
\section{Pilot results and limitations}
At equal duration $T=0.10$ model-time units and fixed smoothing $\sigma=0.50$,
$N=24$ (4 steps), $N=24$ (8 steps), and $N=32$ (4 steps) retain the Alexander
determinant $3\to3$. The $N=24$ relative finite-difference Euler RHS check
changes from $4.28\times10^{-5}$ to $1.07\times10^{-5}$ as $\Delta t$ halves.
This is compatible with second-order central temporal differentiation; it does
not establish long-time stability or an Euler PDE certification.
The same-source N300 geometry is a closed ordered point curve. Separately,
\texttt{CANDIDATE\_SURFACE.npz} is a triangular mesh (3906 vertices, 7808 faces,
$\chi=2$, zero boundary edges), not an ordered vortex line.
No SI core radius, fluid density or other SST canonical constant has been used
to fit a result. The geometry and pilot have no physical unit calibration.
\section{Qualification gates}
\begin{itemize}
\item G09: Knot and link signatures with geometric convergence guards.
\item G10: Kelvin conservation per independently identified material line.
\item G11: Three-dimensional vorticity Euler residual on qualified grid fields.
\item G12: Nonideal PDE reconnection mechanism and convergence, otherwise unresolved.
\item G13: Spatial and temporal convergence at fixed duration and smoothing.
\item G14: Source groups that are truly independent, not differently named copies.
\item G15: Independent SI dimensional calibration of length, time and circulation.
\end{itemize}
\noindent The original DS01--DS03 AES-GCM labels remain sealed. The native
SST Falsifier Framework v1.0.7 is frozen externally and not modified or
impersonated by the present standalone scientific module.
\begin{thebibliography}{9}
\bibitem{Bucher2026}
T. Bucher, A. Gorlach et al. (2026),
``Superluminal correlations in ensembles of optical phase singularities,''
\emph{Nature} \textbf{651}, 920--926.
\url{https://doi.org/10.1038/s41586-026-10209-z}.
\bibitem{Moffatt1969}
H. K. Moffatt (1969), ``The degree of knottedness of tangled vortex lines,''
\emph{Journal of Fluid Mechanics} \textbf{35}, 117--129.
\url{https://doi.org/10.1017/S0022112069000991}.
\bibitem{Saffman1992}
P. G. Saffman (1992), \emph{Vortex Dynamics}, Cambridge University Press.
\url{https://doi.org/10.1017/CBO9780511624063}.
\end{thebibliography}
\end{document}
